% !TEX root = ../algebra_02.tex

\section{Минорный ранг}

\begin{Definition}{Минор}{}
	\emph{Минором} порядка $s$ матрицы $A \in k^{m \times n}$ называется определитель подматрицы $A[\{i_1,\ldots,i_s\},\{j_1,\ldots,j_s\}]$, образованной $s$ строками и $s$ столбцами, стоящими на их пересечении.
\end{Definition}

\begin{Definition}{Минорный ранг}{}
	\emph{Минорным рангом} матрицы $A$ называется максимальный порядок ненулевого минора:
	\[
	\mathrm{rk}_{\mathrm{min}}(A) = \max\{s \mid \exists \text{ ненулевой минор порядка } s\}.
	\]
\end{Definition}

\begin{Theorem}{Все ранги совпадают}{}
	Для любой матрицы $A \in k^{m \times n}$:
	\[
	\mathrm{rk}_{\mathrm{col}}(A) = \mathrm{rk}_{\mathrm{row}}(A) = \mathrm{rk}_{\mathrm{min}}(A).
	\]
\end{Theorem}

\begin{proof}
	Обозначим через $r = \mathrm{rk}(A)$ общее значение столбцового и строкового рангов (которые равны). Покажем два неравенства для минорного ранга.
	
	\textbf{$\mathrm{rk}_{\mathrm{min}}(A) \leq r$.} 
	Рассмотрим произвольный минор порядка $s > r$. Он является определителем некоторой квадратной подматрицы $C$ размера $s \times s$. Так как $\mathrm{rk}_{\mathrm{col}}(A) = r < s$, любые $s$ столбцов матрицы $A$ линейно зависимы. Значит, определитель подматрицы $C$ равен нулю. Следовательно, не существует ненулевых миноров порядка больше $r$.
	
	\textbf{$\mathrm{rk}_{\mathrm{min}}(A) \geq r$.} 
	Поскольку столбцовый ранг равен $r$, существуют $r$ линейно независимых столбцов исходной матрицы с индексами $j_1, \ldots, j_r$. Рассмотрим подматрицу \\ $B = A[{:},\{j_1,\ldots,j_r\}]$ размера $m \times r$. 
	Её столбцовый ранг равен $r$. По теореме о равенстве рангов, строковый ранг подматрицы $B$ также равен $r$.
	Выделив эти строки в матрице $B$, мы получим подматрицу $C = A[\{i_1,\ldots,i_r\},\{j_1,\ldots,j_r\}]$. Эта подматрица квадратная, и её строки линейно независимы. Значит, её определитель отличен от нуля, и мы нашли ненулевой минор порядка $r$.
\end{proof}
